\section{Human Evaluation Protocol}

The classifier-based results above are a preliminary proxy.
We describe here the human evaluation protocol that provides the
ground-truth endpoint for the full study.

\noindent\textbf{Study E1 — Validity audit.}
Participants are shown an original/edited image pair and answer
four questions: (i)~does this appear to be the same place with
a small change; (ii)~does the change appear localised (not global);
(iii)~does the edited image look realistic; and
(iv)~does the change look like a plausible instance of the intended
intervention. Participants respond on a 5-point Likert scale. An edit
is \emph{valid} if it achieves mean score ${\geq}4$ on all four items.

\noindent\textbf{Study E2 — Directional preference.}
Participants complete a 2-alternative forced choice (2AFC) task:
``Which image looks safer?'' Pairs are presented in randomised
left/right order. Five raters per pair. The pair-level
preference share $p_i$ is the primary quantity used to determine
membership in $E(x,a)$. The 2AFC is run only on edits that pass
Study~E1.

\noindent\textbf{Design note.}
E1 and E2 are run as separate studies to prevent the validity
judgement task from contaminating the safety preference response.
Both are deployed via Prolific with attention checks, response-time
filters, and image-load verification. The threshold $\theta$ for
``effective'' is set at preference share ${\geq}0.65$ (one-sided
binomial test, $\alpha{=}0.05$) across $n{=}5$ raters per pair.

\noindent\textbf{Expected outcomes.}
\todo{[NOTE TO AUTHORS: Run E1 + E2 on the pilot subset before the
journal submission. For April 7, this section describes the protocol
and motivation. Report any available results from E1/E2 if the study
launches in time; otherwise present as planned work.]}
